\documentclass[12pt]{article}
\usepackage[utf8]{inputenc}
\usepackage[T1]{fontenc}
\usepackage{hyperref}
\title{Wae Rana Language Resources for K-12 Curriculum Development}
\author{Compiled by Claude Sonnet 5.5 (claude-sonnet-5-5) for Andr\'as Kornai}
\date{}
\begin{document}
\maketitle

\section{Introduction}
Wae Rana (ISO 639-3 wrx; Glottocode waer1237; also called Kolor) is an Austronesian language of central Flores, Indonesia. Wikipedia gives about 3,000 native speakers as of 2010. Sources disagree on its subgrouping. Wikipedia places it in the Central Flores branch of the Ende--Manggarai group, while Verbix lists it under Bima-Lembata and rates it threatened. Joshua Project says the Wae Rana live in two enclaves in south-central Flores, that Wae Rana remains the main language of daily communication, and that Indonesian is widely used for education and trade.

\section{Primary Educational Resources}
The searches found no primers, textbooks or school programmes in Wae Rana.

\section{Linguistic and Reference Materials}
\begin{itemize}
\item Wikipedia, Wae Rana language: \url{https://en.wikipedia.org/wiki/Wae_Rana_language}
\item Verbix, Endangered Languages: Wae Rana: \url{https://docs.verbix.com/EndangeredLanguages/WaeRana}
\end{itemize}
No dictionary or grammar was found in the searches.

\section{Religious and Community Resources}
\begin{itemize}
\item Joshua Project, Wae Rana people group: \url{https://joshuaproject.net/people_groups/19303}
\end{itemize}

\section{Digital and Audio Resources}
None were found.

\section{Curriculum Development Context}
The language is still used in daily life, but I found no written materials. Indonesian is the language of schooling. A curriculum would need new orthography and documentation work with the community, perhaps as local-content (muatan lokal) lessons.
\end{document}